% FeatureScalpel research paper: working template (IEEE conference format).
% Compile: tectonic -X compile main.tex   (or Overleaf, pdfLaTeX + BibTeX).
% Lines starting with % are guidance for you; delete them as you fill each part.
\documentclass[conference]{IEEEtran}
\IEEEoverridecommandlockouts
\usepackage{cite}
\usepackage{amsmath,amssymb,amsfonts}
\usepackage{graphicx}
\usepackage{booktabs}
\usepackage{xcolor}
\usepackage{url}
\usepackage{hyperref}

\newcommand{\todo}[1]{{\color{red}[TODO: #1]}}

\begin{document}

\title{Title of the Paper}

\author{
\IEEEauthorblockN{Author One, Author Two, Author Three, Author Four}
\IEEEauthorblockA{\textit{Department of Computer Science and Engineering (AI)} \
\textit{Mar Baselios College of Engineering and Technology}\
Thiruvananthapuram, India \
email@example.com}
}

\maketitle

% ------------------------------------------------------------------ ABSTRACT
% About 150 to 200 words, one paragraph, no citations, no equations.
% Order: problem, what we did, main numbers, what it means, limits.
\begin{abstract}
\todo{write the abstract last}
\end{abstract}

\begin{IEEEkeywords}
% four to six terms, separated by commas
\todo{keywords}
\end{IEEEkeywords}

% ------------------------------------------------------------------ 1 INTRODUCTION
% Roughly 1 to 1.5 columns. Cite: the problem, the existing fixes, the idea you build on.
\section{Introduction}\label{sec:intro}
\todo{problem, gap, contribution list}

% ------------------------------------------------------------------ 2 RELATED WORK
% Roughly 1 column. Group by topic, not by paper. Cite every claim about other work.
\section{Related Work}\label{sec:related}
\todo{knowledge editing; SAE steering; causal localisation; model stitching}

% ------------------------------------------------------------------ 3 METHOD
% Roughly 2 columns. Cite: methods you use but did not invent (SAE, Adam, ROME's CounterFact).
\section{Method}\label{sec:method}
\subsection{Setting and Notation}
\subsection{Transient Edit through SAE Features}
\subsection{Choosing the Edit}

% ------------------------------------------------------------------ 4 EXPERIMENTS
% Roughly 2 columns. Cite: datasets, metrics, statistical tests. Own results need no citation.
\section{Experimental Setup}\label{sec:setup}
\todo{benchmark, splits, metrics, baselines, controls, hardware}

% ------------------------------------------------------------------ 5 RESULTS
% Roughly 2 columns. One table or figure per claim. Say what each shows and what it does not.
\section{Results}\label{sec:results}

% Figure example (use \columnwidth; figures live in figures/):
% \begin{figure}[t]
% \centering
% \includegraphics[width=\columnwidth]{figures/example.png}
% \caption{What is plotted, on how many items, and the key numbers.}
% \label{fig:example}
% \end{figure}

% Table example (booktabs, no vertical lines):
% \begin{table}[t]
% \caption{Caption goes above a table.}
% \label{tab:example}
% \centering
% \begin{tabular}{lcc}
% \toprule
% Method & Rank 1 (\%) & KL \
% \midrule
% A & 0.0 & 0.0 \
% \bottomrule
% \end{tabular}
% \end{table}

% ------------------------------------------------------------------ 6 DISCUSSION
% Roughly 0.5 to 1 column. What the results support, what they do not, threats to validity.
\section{Discussion and Limitations}\label{sec:discussion}

% ------------------------------------------------------------------ 7 CONCLUSION
% Roughly 0.3 column. No new results, no citations needed.
\section{Conclusion}\label{sec:conclusion}

% Optional, before the references:
% \section*{Acknowledgment}

% ------------------------------------------------------------------ REFERENCES
% An empty reference list stops the compile, so the example entry is cited here.
% Delete the next line as soon as you cite your first real paper.
\nocite{example2000key}
\bibliographystyle{IEEEtran}
\bibliography{references}

\end{document}
